\section{PROPRIET\`A MAGNETICHE DELLA MATERIA}

\subsection{Campi variabili nella materia}

Nel capitolo precedente abbiamo ricavato le equazioni per il campo elettrico in
presenza di materia nel caso statico. Vogliamo stabilire come si modificano le
equazioni di Maxwell se le sorgenti macroscopiche (le cariche e le correnti)
dipendono dal tempo.\\
Il problema \`e sempre quello di passare dalle (14.4), valide per i campi
microscopici, alle equazioni per i campi macroscopici $\vec E$ e $\vec B$, medie
spaziali dei campi microscopici:
\begin{equation}
  \left\{\begin{aligned}
    &\dive\vec E=4\pi(\rho+\rho')\\[4pt]
    &\rot\vec E+\frac{1}{c}\frac{\partial\vec B}{\partial t}=0\\[4pt]
    &\dive\vec B=0\\[4pt]
    &\rot\vec B=\frac{4\pi}{c}(\vec J+\vec J\,')
    +\frac{1}{c}\frac{\partial\vec E}{\partial t}\,\cdot
  \end{aligned}\right.
  \tag{15.1}
\end{equation}
$\vec J\,'$ \`e la media, analoga a $\rho'$, delle correnti microscopiche
$\vec J_\mathrm{micr}$ dovute al moto delle cariche atomiche e molecolari. Nelle
(15.1) le sorgenti $\rho$, $\rho'$, $\vec J$, $\vec J\,'$ possono dipendere
anche dal tempo.

La polarizzazione della materia \`e ancora completamente descritta dal vettore
$\vec P(\vec r,t)$ e le cariche di polarizzazione $\rho'$ e $\sigma'$ sono date
dalle (14.22).\\
Nell'ipotesi di comportamento lineare del mezzo, $\vec P$ dipende linearmente da
$\vec E$ e per campi lentamente variabili con buona approssimazione \`e ancora
dato dalla (14.25), ma in generale la relazione \`e pi\`u complicata
\big($\vec P(\vec r,t)=\int\chi(t-t')\vec E(\vec r,t')\,dt'$\,\big): si pensi,
per esempio, al caso in cui il campo elettrico varia nel tempo in moto
sinusoidale, $\vec E(\vec r,t)=\vec E_0(\vec r\,)\sen\omega t$ (onda
monocromatica), e a un mezzo non polare; sappiamo allora che
nell'approssimazione lineare gli atomi o le molecole sono equivalenti a
oscillatori armonici (l'approssimazione delle piccole oscillazioni). In presenza
di un campo oscillante, dalla teoria degli oscillatori forzati (Volume 1
capitolo 11) sappiamo che l'ampiezza di oscillazione, in questo caso del momento
di dipolo indotto, \`e s\`i proporzionale al campo esterno, ma il coefficiente
di proporzionalit\`a, cio\`e la polarizzabilit\`a $\alpha$, \`e funzione di
$\omega$. Inoltre il dipolo oscilla sfasato rispetto al campo, e ci\`o d\`a
luogo ad un assorbimento di energia da parte dell'oscillatore.\\
Questi effetti diventano rilevanti solo quando ci si avvicina alla frequenza di
risonanza dei nostri oscillatori, che \`e la frequenza $\nu_0$ delle
oscillazioni proprie del sistema; siccome la polarizzazione degli atomi e delle
molecole non polari \`e dovuta alla deformazione della nube elettronica, la
frequenza di risonanza si trova nella regione in cui il sistema emette
radiazione per effetto del moto degli elettroni (transizioni elettroniche),
cio\`e nel campo della luce visibile ($\nu\approx10^{15}\,\mathrm{s}^{-1}$).\\
Il comportamento delle molecole polari in un campo oscillante \`e abbastanza
diverso: ora la polarizzabilit\`a delle molecole non \`e dovuta (solo) agli
elettroni che, come abbiamo visto, riescono a seguire il campo fino a frequenze
piuttosto elevate, ma all'intera molecola, che ha una inerzia ben maggiore di
quella degli elettroni: l'analogo meccanico di
una molecola polare in un campo oscillante \`e un pendolo fisico (per di pi\`u
smorzato) in un campo di gravit\`a $g\sen\omega t$. Il risultato \`e che la
molecola riesce a seguire il campo solo per frequenze relativamente basse
(dipende molto dalla sostanza, ma valori tipici sono intorno ai
$10^6\,\mathrm{s}^{-1}$), al disopra delle quali la polarizzabilit\`a diventa
praticamente nulla.

Consideriamo l'ultima delle (15.1): se prendiamo la divergenza di entrambi i
membri si ottiene
\begin{equation}
  0=\dive\vec J+\dive\vec J\,'
  +\frac{1}{4\pi}\frac{\partial\dive\vec E}{\partial t}
  \tag{15.2}
\end{equation}
e dalla 1$^\mathrm{a}$ delle (15.1) e dall'equazione di conservazione della
carica \ $\dive\vec J=-\partial\rho/\partial t$ \ si ha
\begin{equation}
  \dive\vec J\,'+\frac{\partial\rho'}{\partial t}=0
  \tag{15.3}
\end{equation}
che esprime la conservazione delle cariche atomiche; se utilizziamo la
1$^\mathrm{a}$ delle (14.22) \ $\dive\vec P=-\rho'$ \ si ha
\begin{equation}
  \dive\Big(\vec J\,'-\frac{\partial\vec P}{\partial t}\Big)=0.
  \tag{15.4}
\end{equation}
Dalla (15.4) non segue che \ $\vec J\,'$ deve essere uguale a
\ $\partial\vec P/\partial t$ \ ma, siccome la loro differenza \`e un campo a
divergenza nulla, possono differire per la rotazione di un campo vettoriale
(abbiamo visto nel capitolo 7 che un campo a divergenza nulla pu\`o sempre
essere espresso come la rotazione di un altro campo), quindi poniamo
\begin{equation}
  \vec J\,'=\frac{\partial\vec P}{\partial t}+c\rot\vec M
  \equiv\vec J_p'+\vec J_m'.
  \tag{15.5}
\end{equation}
Il primo termine a 2$^\circ$ membro ha una interpretazione molto semplice, in
quanto rappresenta la corrente media dovuta allo spostamento relativo delle
cariche che danno luogo ai momenti di dipolo atomici o molecolari. Infatti
supponiamo, per semplicit\`a, che i dipoli nel volumetto $\Delta V$ siano
costituiti da cariche puntiformi $q$ e $-q$, tutte uguali (altrimenti
consideriamo i dipoli con cariche comprese fra $q$ e $q+dq$ e poi sommiamo sui
diversi gruppi di dipoli); si ha
\begin{equation}
  \begin{aligned}
    &\vec P(\vec r,t)=\frac{n}{n\Delta V}\sum_{i\in\Delta V}q\vec a_i
    \quad\ \Rightarrow\\
    &\frac{\partial\vec P}{\partial t}=\frac{n}{n\Delta V}
    \sum_{i\in\Delta V}q\dot{\vec a}_i
    =\frac{n}{n\Delta V}\sum_{i\in\Delta V}q(\vec u_i^{\,+}-\vec u_i^{\,-})
  \end{aligned}
  \tag{15.6}
\end{equation}
dove $n$ \`e la densit\`a di dipoli, $\vec u^{\,+}$ e $\vec u^{\,-}$ sono
rispettivamente le velocit\`a delle cariche positive e di quelle negative.\\
Siccome $\rho=nq$ e la velocit\`a media \`e
\ $\vec u\equiv(1/n\Delta V)\sum_i\vec u_i$, si ha
\begin{equation}
  \frac{\partial\vec P}{\partial t}
  =\rho^+\vec u^{\,+}+\rho^-\vec u^{\,-}\equiv\vec J_p'.
  \tag{15.7}
\end{equation}
Per quanto riguarda il termine \ $\vec J_m'\equiv c\rot\vec M$ \ nella (15.5),
prima di tutto dobbiamo convincerci che non solo effettivamente pu\`o essere
diverso da zero, poi vedremo come
\`e legato alle propriet\`a microscopiche della materia.\\
Per quanto riguarda la sua esistenza, tutti sappiamo (e lo sapevano anche gli
antichi) che esistono i magneti, i quali generano nello spazio un campo
magnetico in assenza di sorgenti esterne come correnti macroscopiche $\vec J$ o
campi elettrici: in questi casi l'unico termine che pu\`o essere non nullo a
secondo membro dell'ultima delle (15.1) \`e $\vec J_m'$, e anzi deve essere non
nullo, altrimenti \ $\dive\vec B=0$, $\rot\vec B=0\ \Rightarrow\ \vec B=0$.\\
Quindi $\vec M$ \`e legato al fenomeno della magnetizzazione e questo intanto
giustifica il nome di \textbf{densit\`a di magnetizzazione} per il campo
$\vec M$, e di \textbf{corrente di magnetizzazione} per $\vec J_m'$. Nei
prossimi paragrafi cercheremo di capirne l'origine fisica.

Possiamo riscrivere le (15.1) tenendo conto della definizione del vettore
induzione dielettrica \ $\vec D\equiv\vec E+4\pi\vec P$ \ e definendo
analogamente il campo $\vec H$:
\begin{equation}
  \vec H\defeq\vec B-4\pi\vec M
  \tag{15.8}
\end{equation}
per cui si ha
\begin{equation}
  \left\{\begin{aligned}
    &\dive\vec D=4\pi\rho\\[4pt]
    &\rot\vec E+\frac{1}{c}\frac{\partial\vec B}{\partial t}=0\\[4pt]
    &\dive\vec B=0\\[4pt]
    &\rot\vec H=\frac{4\pi}{c}\vec J
    +\frac{1}{c}\frac{\partial\vec D}{\partial t}\,\cdot
  \end{aligned}\right.
  \tag{15.9}
\end{equation}
Le (15.9), che non presuppongono il comportamento lineare del mezzo, caso per
caso devono essere completate dalle relazioni fra $\vec D$ ed $\vec E$ e fra
$\vec H$ e $\vec B$, dette \textbf{equazioni costitutive}: \`e appunto nelle
equazioni costitutive che si legge se il comportamento del mezzo \`e lineare o
no.

La terminologia relativa ai campi $\vec B$ e $\vec H$ \`e cambiata rispetto ad
alcuni decenni fa: oggi, quasi universalmente, quando si dice ``campo
magnetico'' si intende $\vec B$, e ad $\vec H$ non viene dato nessun nome,
mentre una volta era $\vec H$ che veniva chiamato ``campo magnetico'', e
\ $\vec B=\vec H+4\pi\vec M$ \ ``induzione magnetica''; il motivo va ricercato
nel fatto che il campo generato dai magneti era noto e studiato prima di quello
generato dalle correnti, e quindi era considerato il campo fondamentale. Oggi si
preferisce abbandonare la vecchia terminologia per mettere in risalto il fatto
che il campo fondamentale \`e $\vec B$, in quanto \`e il campo che interviene
nelle equazioni microscopiche (14.4), e $\vec H$ nasce solo in seguito al
processo di media che porta alla descrizione macroscopica: i campi ausiliari
$\vec D$ e $\vec H$ potrebbero non venire mai introdotti, in quanto sarebbero
sufficienti i campi $\vec E$ e $\vec B$, pi\`u \ $\vec P$ e $\vec M$ \ per
descrivere le propriet\`a elettriche e magnetiche della materia.

Nel sistema di Gauss $\vec B$ ed $\vec H$ sono grandezze omogenee, cio\`e non
solo hanno le stesse dimensioni, ma possono anche essere sommate
\big(v.\ (15.8)\big), e quindi l'unit\`a di misura \`e la stessa, cio\`e il
gauss; tuttavia, come eredit\`a di un passato ancora non abbastanza lontano
resta la consuetudine di chiamare \textbf{oersted} (abbreviato Oe) questa stessa
unit\`a di misura (il gauss), quando ci si riferisce al campo $\vec H$: cos\`i,
se siamo nel vuoto dove $\vec B=\vec H$, un campo $B$ di $10^3$\,gauss, \`e un
campo $H$ di $10^3$\,oersted!

\subsection{Fenomenologia}

Il comportamento della materia in un campo magnetico ha qualche somiglianza con
quello dei dielettrici in un campo elettrico, ma \`e pi\`u articolato. Mentre
abbiamo visto che un dielettrico in un campo elettrico non omogeneo viene sempre
attirato nella direzione in cui aumenta l'intensit\`a del campo, con una forza
che \`e proporzionale al quadrato del campo (se il campo viene raddoppiato la
forza aumenta di un fattore 4), nel caso di sostanze poste un campo magnetico
non omogeneo, come ad esempio vicino agli estremi di un solenoide, alcune
vengono attratte dal solenoide, cio\`e verso la regione in cui il campo \`e
pi\`u intenso, altre vengono respinte.\\
Inoltre, nel caso delle sostanze che vengono attratte dal solenoide, per la
maggior parte di esse la forza \`e, come nel caso del campo elettrico,
proporzionale al quadrato del campo magnetico ed \`e sempre molto debole (saremo
pi\`u precisi pi\`u avanti) e decresce all'aumentare della temperatura; le
sostanze di questo tipo sono dette \textbf{paramagnetiche}: ad esempio, sono
paramagnetiche l'alluminio, il sodio, il platino, il cloruro di rame (CuCl$_2$),
ecc.\\
Alcune sostanze, invece, vengono attratte nella direzione in cui il campo
aumenta con una forza molto intensa proporzionale al campo, e non al suo
quadrato: queste sostanze sono dette \textbf{ferromagnetiche}, e rientrano in
questa classe il ferro, il cobalto, il nichel, purch\'e al disotto di una
opportuna temperatura, detta \textbf{temperatura di Curie} $\Theta$, che dipende
dalla sostanza: per il ferro $\Theta=774\,^\circ$C, per il cobalto
$\Theta=1131\,^\circ$C, e per il nichel $\Theta=372\,^\circ$C; al disopra della
temperatura di Curie le sostanze ferromagnetiche diventano paramagnetiche.\\
Infine, nel caso delle sostanze che vengono respinte dal solenoide, la forza \`e
molto debole (normalmente circa un ordine di grandezza pi\`u piccola di quella
che agisce sulle sostanze paramagnetiche a temperatura ambiente), e sempre
proporzionale al quadrato del campo, ed inoltre questa forza \`e indipendente
dalla temperatura: queste sostanze sono dette \textbf{diamagnetiche}, e fra esse
troviamo l'H$_2$O, il rame, il piombo, l'oro, lo zolfo e il cloruro di sodio.

Se lasciamo da parte (per ora) le sostanze ferromagnetiche e facciamo il
confronto con il caso dei dielettrici, in primo luogo notiamo che tutti i
dielettrici vengono attirati nella direzione in cui aumenta l'intensit\`a del
campo, mentre nel caso magnetico questo \`e vero solo per le sostanze
paramagnetiche (quindi forse sarebbe pi\`u appropriato chiamare
``paraelettrici'' i dielettrici \dots\ ): il comportamento dei dielettrici \`e
dovuto al fatto che la costante dielettrica \`e sempre maggiore di 1 (o
equivalentemente la suscettivit\`a elettrica $\chi_e>0$), e questo perch\'e sia
il momento di dipolo indotto che il momento di dipolo medio (nel caso delle
sostanze polari) \`e sempre nello stesso verso del campo.\\
In secondo luogo, sempre nel caso dei dielettrici, la dipendenza delle
propriet\`a dielettriche dalla temperatura si ha solo per le sostanze polari,
cio\`e quelle in cui la polarizzazione \`e dovuta all'orientamento dei momenti
di dipolo permanenti (o propri), mentre \`e assente per le sostanze non polari;
nel caso magnetico solo le sostanze paramagnetiche hanno un comportamento che
dipende dalla temperatura.\\
Infine, notiamo che le espressioni della forza e del momento che un campo
magnetico esercita su un momento magnetico sono formalmente identiche a quelle
della forza e del momento che un campo elettrico esercita su un dipolo, pur di
sostituire $\vec p$ con $\vec m$ ed
$\vec E$ con $\vec B$: i momenti sono dati rispettivamente dalla (7.59) e dalla
(14.15); per quanto riguarda le forze, tenuto conto che (nel caso statico) il
campo elettrico \`e conservativo, e quindi \ $\partial_iE_j=\partial_jE_i$, la
(14.62) pu\`o essere riscritta come \big(v.\ anche (14.65)\big)
\begin{equation}
  F_i=\textstyle\sum_jp_j\partial_jE_i=\sum_jp_j\partial_iE_j
  \quad\ \Rightarrow\quad\ \vec F=\nabla(\vec p\cdot\vec E)
  \tag{15.10}
\end{equation}
e la (7.62), come abbiamo gi\`a detto a suo tempo, in generale si scrive
\begin{equation}
  \vec F=\nabla(\vec m\cdot\vec B).
  \tag{15.11}
\end{equation}
Questo confronto suggerisce che nel passaggio dalla trattazione dei dielettrici
a quella delle sostanze paramagnetiche e diamagnetiche il momento magnetico
$\vec m$ degli atomi o delle molecole prenda il posto del momento di dipolo
$\vec p$; la densit\`a di magnetizzazione, definita come
\begin{equation}
  \vec M(x,y,z)=\frac{\sum_i\vec m_i}{\Delta V}
  \tag{15.12}
\end{equation}
quello della densit\`a di polarizzazione $\vec P$; che nelle sostanze
diamagnetiche il momento magnetico sia un momento magnetico \underline{indotto},
in verso \underline{opposto} al campo magnetico inducente, mentre le sostanze
paramagnetiche siano quelle i cui atomi possiedono un momento magnetico proprio
(o permanente), che parzialmente si allinea con il campo.\\
Inoltre, nell'ipotesi che il comportamento magnetico del mezzo sia lineare
(ipotesi ampiamente soddisfatta da tutte le sostanze paramagnetiche e
diamagnetiche), $\vec M$ \`e proporzionale a $\vec B$, e la costante di
proporzionalit\`a $\chi_m$ \`e detta \textbf{suscettivit\`a magnetica}; in
effetti, sempre come eredit\`a dei tempi in cui $\vec H$ era considerato il
campo fondamentale, anzich\'e \ $\vec M=\chi_m\vec B$ \ si pone
\begin{equation}
  \vec M=\chi_m\vec H.
  \tag{15.13}
\end{equation}
In pratica le due definizioni, la (15.13) e quella pi\`u naturale ma non
adottata \ $\vec M=\chi_m\vec B$, sono equivalenti perch\'e, grazie alla (15.8),
differiscono per termini di ordine $(\chi_m)^2$, e $\chi_m$ \`e sempre molto
piccolo: \ $|\chi_m|\simeq10^{-6}$ \ per le sostanze diamagnetiche, ed arriva
eccezionalmente ad assumere il valore di $\simeq10^{-3}$ nel caso dell'ossigeno
liquido (paramagnetico), ma in questo caso (ovviamente) la temperatura \`e molto
bassa.\\
La suscettivit\`a magnetica $\chi_m$ dovrebbe essere (ed \`e) negativa per le
sostanze diamagnetiche, e positiva per quelle paramagnetiche.\\
In tutti i casi in cui $\vec M$ \`e proporzionale ad $\vec H$ (o
equivalentemente a $\vec B$) si definisce la \textbf{permeabilit\`a magnetica}
$\mu$, analoga alla costante dielettrica $\epsilon$, come
\begin{equation}
  \mu\defeq1+4\pi\chi_m
  \tag{15.14}
\end{equation}
per cui in questi casi dalla (15.8) e (15.13) si ha
\begin{equation}
  \vec B=\mu\vec H.
  \tag{15.15}
\end{equation}
La permeabilit\`a magnetica $\mu$ pu\`o essere maggiore o minore di 1 e, come
abbiamo visto, per le sostanze non ferromagnetiche differisce poco da 1, per cui
gli effetti magnetici sono generalmente trascurabili.

\subsection{Il momento magnetico degli atomi}

Nel paragrafo precedente abbiamo illustrato il comportamento della materia in un
campo magnetico e per analogia con la trattazione dei dielettrici abbiamo
anticipato quello che dovrebbe essere il risultato di una trattazione teorica
del problema. Il punto \`e che, a stretto rigore, non esiste una teoria classica
del magnetismo nella materia, dato che secondo la fisica classica (o meglio, la
meccanica statistica classica) non esiste n\'e il diamagnetismo, n\'e il
paramagnetismo, n\'e il ferromagnetismo: sono tutti fenomeni che possono essere
correttamente spiegati solo facendo ricorso alla meccanica quantistica, in
quanto intervengono in maniera determinante sia lo spin dell'elettrone, cio\`e
il suo ``momento angolare intrinseco'', sia il principio di esclusione di Pauli,
sia gli effetti legati alla quantizzazione.\\
Quindi non potremo andare molto lontano con una trattazione puramente classica e
il linguaggio classico che useremo sar\`a spesso inadeguato.

Nel capitolo 7.5 abbiamo definito il momento magnetico $\vec m$ di una spira
percorsa da corrente; in generale, se abbiamo una distribuzione di volume
$\vec J$ di corrente contenuta in un volume finito $V$, il momento magnetico
$\vec m$ \`e definito come
\begin{equation}
  \vec m\defeq\frac{1}{2c}\int_V\vec r\,'\wedge\vec J(\vec r\,')\,dV.
  \tag{15.16}
\end{equation}
Innanzi tutto la (15.16) coincide con la (7.49) se la corrente \`e stazionaria
ed \`e concentrata in una spira \big(brevemente: si usa l'identit\`a
\ $\hat n\cdot(\vec r\wedge d\vec l\,)=-(\vec r\wedge\hat n)\cdot d\vec l$
\ dove $\hat n$ \`e un versore arbitrario, poi si usa il teorema di Stokes e
l'identit\`a \ $\rot(\vec r\wedge\hat n)=-2\hat n$\,\big); in secondo luogo,
sempre nel caso di una corrente stazionaria ($\dive\vec J=0$), con la
definizione (15.16) vale la (7.50), cio\`e
\[
  \vec A(\vec r\,)=\frac{\vec m\wedge\vec r}{r^3}+\cdots.
\]
Consideriamo ora un atomo; gli elettroni che ruotano attorno al nucleo danno
luogo ad una distribuzione di corrente rapidamente variabile nel tempo; se
per\`o facciamo la media di questa corrente su un tempo grande rispetto ai
periodi orbitali degli elettroni ($T\approx10^{-15}$\,s), sparisce la dipendenza
dal tempo e il momento magnetico (medio) dell'atomo dato dalla (15.16) \`e
costante. (La media temporale trasforma, per esempio, un elettrone che ruota su
un'orbita circolare di raggio $a$ con velocit\`a $v$ in una spira percorsa dalla
corrente stazionaria $I=(e/2\pi a)v$; in una trattazione quantistica la media
temporale non \`e necessaria perch\'e la corrente in un atomo \`e
automaticamente stazionaria).\\
Il processo di media che portato alla distribuzione di corrente
\ $\vec J=\rho\vec v$, dove $\rho$ \`e la distribuzione di carica degli
elettroni e $\vec v$ \`e il campo delle velocit\`a di questa distribuzione,
porta anche ad una distribuzione di massa $\mu$ degli elettroni proporzionale a
$\rho$ : \ $\mu=-(m/e)\rho$ \ ($m$ \`e la massa dell'elettrone), e quindi
\ $\vec J=-(e/m)\mu\vec v$. Se inseriamo questa relazione nella (15.16) si ha
\begin{equation}
  \vec m=-\frac{e}{2mc}\int_V(\vec r\,'\wedge\vec v\,)\mu(\vec r\,')\,dV.
  \tag{15.17}
\end{equation}
L'integrale nella (15.17) \`e il \textit{momento angolare} $\vec L$ dell'atomo
che, in seguito al processo di media, viene ora considerato come un sistema
continuo; quindi
\begin{equation}
  \vec m=-\frac{e}{2mc}\vec L.
  \tag{15.18}
\end{equation}
Il momento magnetico dato dalla (15.18) \`e detto ``momento magnetico
orbitale'', in quanto \`e dovuto al \textit{moto orbitale} degli elettroni.\\
Inoltre, \textit{ogni} elettrone ha un momento magnetico intrinseco (cio\`e
anche quando \`e fermo), detto anche momento magnetico di spin, anch'esso
proporzionale al momento angolare intrinseco $\vec s$, ma con coefficiente di
proporzionalit\`a $-e/mc$, cio\`e il doppio di quello che compare nella
(15.18).\\
Il momento magnetico intrinseco, come il momento angolare di spin, \`e di natura
completamente quantistica: per aiutare l'intuizione uno pu\`o pensare
all'elettrone come se fosse una sferetta carica rotante; questo per\`o \`e solo
un modello che non \`e in grado di spiegare il comportamento dell'elettrone.\\
Il momento magnetico dell'atomo \`e la somma (vettoriale) dei \textit{momenti
magnetici orbitali e di quelli di spin} e se non \`e nullo allora diciamo che
l'atomo ha un momento magnetico proprio. Secondo la meccanica quantistica il
momento angolare \`e ``quantizzato'', cio\`e i valori che pu\`o assumere sono
solo i multipli interi della costante di Planck, divisa per $2\pi$ :
\ $\hbar\equiv h/2\pi\approx10^{-27}\,$erg\,s, e siccome il momento angolare di
un atomo, quando non \`e nullo, \`e sempre di poche unit\`a di $\hbar$, l'ordine
di grandezza del momento magnetico proprio di un atomo \`e
\begin{equation}
  m_{\textsc b}=\frac{e\hbar}{2mc}=\frac{4.8\times10^{-10}\cdot10^{-27}}
  {2\cdot0.9\times10^{-27}\cdot3\times10^{10}}
  \simeq0.9\times10^{-20}\,\mathrm{erg/gauss}
  \tag{15.19}
\end{equation}
$m_{\textsc b}$ \`e una costante fondamentale ed \`e chiamata \textit{magnetone
di Bohr}.

In un campo magnetico le sostanze i cui atomi possiedono un momento magnetico
proprio si comportano come le sostanze polari in un campo elettrico, quindi non
ci soffermiamo ulteriormente su di esse, ma preferiamo discutere il caso delle
sostanze i cui atomi non hanno un momento magnetico permanente in quanto, come
abbiamo gi\`a anticipato, il loro comportamento in un campo esterno presenta
differenze rispetto a quello delle molecole non polari in un campo elettrico.\\
Se un atomo non possiede un momento magnetico proprio, un campo magnetico
esterno induce nell'atomo un momento magnetico secondo il meccanismo che ora
descriviamo.

Cominciamo con un caso semplice: immaginiamo di avere un elettrone che ruota
attorno ad un nucleo di carica $Ze$ descrivendo un'orbita circolare di raggio
$r$ con velocit\`a angolare $\omega_0$ data da
\begin{equation}
  m\omega_0^2r=\frac{Ze^2}{r^2}\,\cdot
  \tag{15.20}
\end{equation}
Mettiamo l'atomo dentro un solenoide con l'asse ortogonale all'orbita
dell'elettrone e supponiamo che inizialmente $\vec B=0$ $(I=0)$; quindi
accendiamo la corrente nel solenoide finch\'e il campo raggiunge un certo valore
$\vec B\neq0$. Poich\'e il campo magnetico varia, si ha un campo elettrico
indotto, e siccome la sua circuitazione lungo l'orbita dell'elettrone non \`e
nulla, il suo effetto \`e quello di modificare la velocit\`a orbitale
dell'elettrone: per la
legge di Lenz, siccome $\dot B>0$, la variazione della velocit\`a dell'elettrone
\`e in verso tale da dar luogo ad un campo magnetico opposto a quello del
solenoide quindi, siccome la carica dell'elettrone \`e negativa, se l'elettrone
ruota in verso orario rispetto al verso di $\vec B$ $(\omega_0<0)$ viene
rallentato, altrimenti $(\omega_0>0)$ viene accelerato e alla fine la velocit\`a
angolare $\omega$ dell'elettrone sar\`a diversa da quella iniziale $\omega_0$.\\
Per calcolare $\omega$ supponiamo che sia trascurabile la variazione del raggio
dell'orbita, cio\`e che la forza centrale $Ze^2/r^2$ sia molto maggiore della
forza di Lorentz; allora la componente centripeta di $\vec F=m\vec a$, tenuto
conto che ora c'\`e anche la forza di Lorentz, \`e
\begin{equation}
  m\omega^2r=\frac{Ze^2}{r^2}+\frac{eB}{c}\omega r.
  \tag{15.21}
\end{equation}
Posto
\begin{equation}
  \omega_{\textsc l}=\frac{eB}{2mc}
  \tag{15.22}
\end{equation}
e tenuto conto della (15.20) si ha
\begin{equation}
  \omega^2-2\omega\omega_{\textsc l}-\omega_0^2=0\quad\Rightarrow
  \quad\omega=\omega_{\textsc l}\pm\sqrt{\omega_0^2+\omega_{\textsc l}^2}
  \tag{15.23}
\end{equation}
e se trascuriamo $\omega_{\textsc l}^2$ rispetto a $\omega_0^2$ si ha, sia per
$\omega_0>0$ che per $\omega_0<0$,
\begin{equation}
  \omega=\omega_0+\omega_{\textsc l}.
  \tag{15.24}
\end{equation}
Prima di tutto verifichiamo in quali condizioni \`e legittimo trascurare
$\omega_{\textsc l}^2$ rispetto a $\omega_0^2$ (mettiamo $Z=1$):
\begin{equation}
  \frac{\omega_{\textsc l}^2}{\omega_0^2}
  =\frac{e^2B^2}{4m^2c^2}\cdot\frac{mr^3}{e^2}=\frac{B^2r^3}{4mc^2};
  \quad\omega_{\textsc l}^2\ll\omega_0^2\ \Rightarrow\ B^2\ll\frac{4mc^2}{r^3}
  \tag{15.25}
\end{equation}
e se poniamo $r=10^{-8}\,$cm si ha che deve essere $B\ll10^{13}\,$gauss, e
questa \`e anche la condizione che deve essere verificata affinch\'e sia
trascurabile la variazione del raggio dell'orbita.\\
Il significato della (15.24) \`e questo: qualunque sia la velocit\`a e il verso
di rotazione dell'elettrone, il campo magnetico induce una variazione della sua
velocit\`a angolare data da $\omega_{\textsc l}$ (\textit{frequenza di Larmor}),
e quindi una variazione $\Delta\vec m$ del momento magnetico associato al suo
moto. Si noti che la frequenza di Larmor \`e la met\`a della velocit\`a angolare
di una carica \underline{libera} in un campo magnetico uniforme, data dalla
(5.11) (con $\gamma(v)=1$): infatti in questo caso $\omega_0=0$ e dalla (15.23)
segue che, se la particella non \`e ferma, $\omega=2\omega_{\textsc l}$.\\
Il risultato espresso dalla (15.24) \`e lo stesso se, invece di accendere il
campo magnetico, l'atomo viene portato da una regione in cui $\vec B=0$ ad una
in cui $\vec B\neq0$: in questo caso \`e la forza di Lorentz dovuta alla
velocit\`a con cui l'atomo viene spostato che produce la variazione della
velocit\`a angolare dell'elettrone.\\
La variazione di $\vec m$ mediata su un periodo \`e data da (l'asse $z$ \`e
orientato come $\vec B$)
\begin{equation}
  \Delta m_z=\frac{\Delta I_{\mathrm{el}}}{c}\,\pi r^2
  =\frac{-e\omega_{\textsc l}r}{2\pi rc}\,\pi r^2=-\frac{e^2r^2}{4mc^2}\,B
  \tag{15.26}
\end{equation}
quindi il momento magnetico indotto \`e \underline{opposto} al campo magnetico:
il segno $-$ viene dalla legge di Lenz, e questa \`e l'origine del
diamagnetismo.

Cos\`i, per esempio, se abbiamo due elettroni che ruotano in versi opposti sulla
stessa orbita attorno ad un nucleo (trascuriamo la repulsione fra gli
elettroni), il momento magnetico proprio dell'atomo \`e nullo, ma se l'atomo
viene posto in un campo magnetico il momento magnetico indotto \`e il doppio di
quello dato dalla (15.26).\\
Se nella (15.26) prendiamo $r=10^{-8}\,$cm, \ $B=10^4\,$gauss si ha
$|\Delta m_z|\simeq10^{-25}\,$erg/gauss, 5 ordini di grandezza inferiore al
magnetone di Bohr: all'incirca lo stesso rapporto che si ha nella teoria dei
dielettrici fra il momento di dipolo indotto da un campo di $10^4\,$V/cm e i
momenti di dipolo permanenti.

Abbiamo capito su un esempio semplice l'origine del diamagnetismo: il risultato
\`e generale, indipendente dal moto degli elettroni all'interno di un atomo, in
quanto il campo magnetico esterno -- tramite il campo elettrico indotto o la
forza di Lorentz -- induce su un sistema di elettroni legati un moto di insieme
di rotazione attorno alla direzione del campo con la frequenza di Larmor, e
quindi un momento magnetico opposto a $\vec B$ (\textbf{teorema di Larmor});
vale la pena di dare una dimostrazione (una delle tante possibili) di questo
fatto.\\
Consideriamo un atomo di numero atomico $Z$ e, per semplicit\`a, supponiamo che
si trovi al centro di un solenoide nel quale viene acceso il campo. Supponiamo
che all'inizio il momento angolare $\vec L$ dell'atomo sia nullo, e quindi per
la (15.18) l'atomo non possiede momento magnetico orbitale proprio. Calcoliamo
mediante il teorema dell'impulso angolare il momento angolare finale dell'atomo.
Le forze che il nucleo esercita su ogni elettrone sono a momento nullo (forze
centrali), e quelle di repulsione fra gli elettroni sono forze interne, quindi a
momento risultante nullo; le uniche forze che hanno un momento non nullo sono le
forze di Lorentz \ $-e\big(\vec E+(\vec v/c)\wedge\vec B\big)$ \ che agiscono
sugli elettroni, dove $\vec E$ \`e il campo elettrico indotto dalla variazione
di $\vec B$ ed \`e dato dalla (8.20). Il momento della forza di Lorentz su un
elettrone \`e
\begin{equation}
  \vec M=-e\vec r\wedge(\vec E+\frac{\vec v}{c}\wedge\vec B);
  \qquad\vec E=\frac{1}{2c}\vec r\wedge\dot{\vec B}.
  \tag{15.27}
\end{equation}
Usiamo l'identit\`a \ $\vec a\wedge(\vec b\wedge\vec c\,)
=(\vec a\cdot\vec c\,)\,\vec b-(\vec a\cdot\vec b\,)\,\vec c$ \ e prendiamo
l'asse $z$ nel verso di $\vec B$; si ha
\begin{equation}
  M_z=-\frac{e}{2c}\Big(z^2\dot B-r^2\dot B+2z\dot zB-2(\vec r\cdot\vec v)B\Big)
  =\frac{e}{2c}\frac{d}{dt}\Big((x^2+y^2)B\Big).
  \tag{15.28}
\end{equation}
Sommando sugli $Z$ elettroni e integrando sul tempo si ottiene
\begin{equation}
  L_z=\frac{eB}{2c}\sum_{i=1}^Z(x^2+y^2)_i
  \equiv Z\frac{eB}{2c}\,(\overline{x^2+y^2}).
  \tag{15.29}
\end{equation}
In modo analogo, ma alquanto pi\`u laborioso, si dimostra che le componenti
$L_x$ e $L_y$ sono (in media) nulle. Quindi, per la (15.18), il momento
magnetico indotto \`e dato da
\begin{equation}
  m_z=-Z\frac{e^2B}{4mc^2}\,(\overline{x^2+y^2}).
  \tag{15.30}
\end{equation}
Siccome sappiamo che l'effetto del campo magnetico sulle dimensioni delle orbite
\`e trascurabile e per un atomo con $\vec L=0$ non c'\`e nessuna direzione
privilegiata,
\ $\overline{x^2}=\overline{y^2}=\overline{z^2}=(1/3)\overline{r^2}$\,, quindi
possiamo scrivere
\begin{equation}
  m_z=-Z\frac{e^2B}{6mc^2}\,\overline{r^2}
  \tag{15.31}
\end{equation}
dove $\overline{r^2}$ \`e il raggio quadratico medio dell'atomo: varia da atomo
ad atomo, ma \`e sempre dell'ordine di $10^{-8}\,$cm.\\
Il problema che abbiamo appena discusso ha molti punti in comune con il problema
della rotazione del piano di oscillazione di un pendolo di Foucault al polo, e
ci\`o \`e dovuto all'identit\`a formale fra la forza di Coriolis in un
riferimento rotante e la forza di Lorentz.\\
(Abbiamo detto che secondo la meccanica quantistica il momento angolare pu\`o
assumere solo valori multipli di $\hbar$, mentre dalla (15.29) risulta che $L_z$
\`e proporzionale a $B$, quindi pu\`o assumere qualsiasi valore: non c'\`e
contraddizione, perch\'e nella formulazione hamiltoniana della meccanica il
momento angolare ``canonico'' \ $\vec L=\vec r\wedge\vec p$ \ in presenza di un
campo magnetico \`e diverso dal momento angolare ``cinetico''
\ $\vec r\wedge(m\vec v\,)$\,, cio\`e $\vec p\neq m\vec v$, ed \`e il momento
angolare canonico che \`e quantizzato, mentre nelle (15.17) e (15.18) interviene
il momento angolare cinetico\big).

Dalla (15.31) segue che il momento magnetico indotto \`e proporzionale al campo
magnetico, e quindi possiamo definire la \textbf{polarizzabilit\`a magnetica}
$\alpha_m$ di un atomo:
\begin{equation}
  \vec m=\alpha_m\vec B\approx\alpha_m\vec H
  \tag{15.32}
\end{equation}
(come abbiamo detto a proposito della (15.13) le definizioni in termini di
$\vec B$ o di $\vec H$ in pratica sono equivalenti), analoga alla
polarizzabilit\`a elettrica definita dalla (14.9). Anche in questo caso la
proporzionalit\`a fra momento magnetico e campo \`e solo approssimata, in quanto
segue dall'aver trascurato nella (15.31) la dipendenza del raggio medio delle
orbite degli elettroni dal campo (come nel modello di atomo che ci ha portato
alla (14.9) abbiamo trascurato la deformazione della distribuzione di carica
negativa), ma la (15.25) ci garantisce che (per le sostanze diamagnetiche) ci\`o
\`e legittimo per i campi che possono essere prodotti in laboratorio.\\
Dalla (15.31) si ha che la polarizzabilit\`a (dia)magnetica $\alpha_m$ di un
atomo con $Z$ elettroni \`e dell'ordine di
\begin{equation}
  \alpha_m=-Z\frac{e^2}{6mc^2}\,\overline{r^2}
  \simeq-Z\cdot5\times10^{-30}\,\mathrm{cm}^3
  \tag{15.33}
\end{equation}
cio\`e parecchi ordini di grandezza $(3\div5)$ inferiore alla polarizzabilit\`a
elettrica ($\alpha_e\approx10^{-24}\,\mathrm{cm}^3$), e ci\`o spiega come mai
gli effetti magnetici nelle sostanze non ferromagnetiche sono quasi sempre
trascurabili.\\
Grazie alla piccolezza di $\alpha_m$ la suscettivit\`a magnetica $\chi_m$
definita dalle (15.12) e (15.13) \`e con ottima approssimazione data da
\begin{equation}
  \chi_m=n\alpha_m=-nZ\frac{e^2}{6mc^2}\,\overline{r^2}
  \tag{15.34}
\end{equation}
($n$ \`e il numero di atomi per unit\`a di volume), in quanto il campo
microscopico prodotto dai momenti magnetici indotti \`e sempre molto piccolo,
per cui non \`e necessario ricercare una formula analoga a quella di
Mossotti--Clausius. Le misure di suscettivit\`a magnetica costituiscono un buon
metodo per ricavare le dimensioni degli atomi delle sostanze diamagnetiche.

\subsection{La corrente di magnetizzazione}

La densit\`a di magnetizzazione $\vec M$ \`e definita dalla (15.12): resta da
vedere che essa coincide con il campo $\vec M$ introdotto in maniera formale
nell'equazione (15.5), cio\`e che effettivamente la densit\`a di magnetizzazione
d\`a luogo a una corrente, detta di magnetizzazione.\\
La dimostrazione formale di questo fatto \`e simile a quella che ci ha permesso
di stabilire le (14.22): si parte dall'identit\`a
\begin{equation}
  M_i=\frac12\Big([\vec r\wedge\rot\vec M]_i
  +\dive(M_i\vec r\,)-\partial_i(\vec r\cdot\vec M)\Big).
  \tag{15.35}
\end{equation}
Consideriamo, per esempio, la componente $x$: se $\hat\imath$ \`e il versore
dell'asse $x$ possiamo scrivere gli ultimi due termini nella (15.35) come
$\dive\big(M_x\vec r-(\vec r\cdot\vec M)\hat\imath\big)$, quindi per il teorema
della divergenza si ha
\begin{equation}
  \begin{split}
    \int_VM_i\,dV&=\frac12\int_V[\vec r\wedge\rot\vec M]_i\,dV
    +\frac12\int_S\big[(\vec r\cdot\hat n)M_i
    -(\vec r\cdot\vec M)\,\hat n_i\big]\,dS\\
    &=\frac12\int_V[\vec r\wedge\rot\vec M]_i\,dV
    +\frac12\int_S\big[\vec r\wedge(\vec M\wedge\hat n)\big]_i\,dS.
  \end{split}
  \tag{15.36}
\end{equation}
D'altronde, se $\vec J_m'$ \`e la corrente di magnetizzazione dentro il volume
$V$ e $\vec{\mathcal F}_m'$ la corrente (di magnetizzazione) sulla superficie
$S$ del volume $V$, per la (15.16) si ha
\begin{equation}
  \int_VM_i\,dV=\frac{1}{2c}\int_V(\vec r\wedge\vec J_m')_i\,dV
  +\frac{1}{2c}\int_S(\vec r\wedge\vec{\mathcal F}_m')_i\,dS.
  \tag{15.37}
\end{equation}
Dal confronto fra la (15.36) e la (15.37) e dal fatto che, scelta opportunamente
l'origine delle coordinate e le dimensioni del volume di integrazione \big(\`e
lo stesso argomento che abbiamo utilizzato per passare dalle (14.20) e (14.21)
alla (14.22)\big), gli integrali di volume sono infinitesimi di ordine maggiore
di quelli di superficie, si conclude l'uguaglianza degli integrandi; ancora
sfruttando l'arbitrariet\`a dell'origine, e quindi del vettore $\vec r$, segue
\begin{equation}
  \vec J_m'=c\rot\vec M;\qquad\vec{\mathcal F}_m'=c\vec M\wedge\hat n
  \tag{15.38}
\end{equation}
Una dimostrazione alternativa consiste nel mostrare che il campo generato da una
magnetizzazione $\vec M$ in un volume arbitrario $V$ genera lo stesso potenziale
vettore delle correnti $\vec J_m'$ e $\vec{\mathcal F}_m'$ date dalla (15.38):
\begin{equation}
  \vec A(\vec r\,)=\int_V\frac{\vec M(\vec r\,')\wedge(\vec r-\vec r\,')}
  {|\vec r-\vec r\,'|^3}\,dV
  =\int_V\frac{\rot\vec M(\vec r\,')}{|\vec r-\vec r\,'|}\,dV
  +\int_S\frac{M(\vec r\,')\wedge\hat n}{|\vec r-\vec r\,'|}\,dS.
  \tag{15.39}
\end{equation}
Il vantaggio della prima dimostrazione \`e che si applica anche al caso in cui
la magnetizzazione dipende dal tempo, mentre la seconda vale nel caso statico.
Va comunque notato che da entrambe le dimostrazioni emerge che la
magnetizzazione in un volume $V$ \`e \underline{equivalente} ad una corrente di
volume $\vec J_m'$ e ad una corrente superficiale $\vec{\mathcal F}_m'$, e non
necessariamente che queste correnti esistono all'interno del volume $V$: infatti
parte della magnetizzazione \`e dovuta al momento magnetico intrinseco degli
elettroni, e non sembra che esso sia attribuibile ad un moto di cariche, cio\`e
a ``correnti amperiane'' (il modello dell'elettrone come una sferetta carica
rotante \`e solo un modello e non sembra corrispondere alla realt\`a).

Per capire l'origine delle (15.38), supponiamo di avere un cilindro
uniformemente magnetizzato parallelamente all'asse: la figura 15.1 illustra una
sezione del cilindro e le correnti atomiche che danno luogo alla magnetizzazione
(si confronti la figura 15.1 con la 14.3). \`E chiaro che la corrente che
attraversa una qualsiasi superficie (molto maggiore delle dimensioni degli
atomi) tutta interna al cilindro \`e nulla: ogni atomo contribuisce in egual
misura ad una corrente che attraversa la superficie in un verso e nel verso
opposto. Invece sulla superficie laterale del cilindro le correnti atomiche
danno luogo ad una corrente superficiale (in verso antiorario nella figura 15.1)
ortogonale alla direzione di $\vec M$. Se invece la magnetizzazione non \`e
uniforme sono possibili anche correnti di volume all'interno del cilindro: per
esempio, se la densit\`a di magnetizzazione varia con la distanza dall'asse, si
hanno correnti di volume circolanti attorno all'asse.
\fig{fig15_1}{15.1}

\subsection{Le sostanze ferromagnetiche}

Supponiamo di avere una calamita di forma cilindrica: siccome una calamita (o
magnete permanente) genera un campo magnetico, necessariamente la
magnetizzazione $\vec M$ deve essere non nulla. Supponiamo che $\vec M$ sia
costante dentro il cilindro, e diretto come l'asse. Vogliamo capire come sono
fatti i campi $\vec B$ ed $\vec H$ all'interno e all'esterno del magnete, se
l'unica sorgente dei campi \`e la magnetizzazione $\vec M$.\\
Il problema pu\`o essere affrontato da due punti di vista diversi, a seconda che
lo si voglia studiare in termini di $\vec B$ o di $\vec H$.

La sorgente di $\vec B$ \`e solo la corrente superficiale di magnetizzazione
che, per la (15.38), \`e data da $\mathcal F=cM$ e circola in verso antiorario
attorno alla direzione di $\vec M$, perci\`o il magnete \`e equivalente ad un
solenoide con le stesse dimensioni del cilindro magnetizzato. Quindi, in
particolare, dentro il cilindro $\vec B$ non \`e uniforme e le sue linee di
forza vanno nello stesso verso di quelle di $\vec M$, sono pi\`u dense al
centro, risultano incurvate verso l'esterno, alcune di esse attraversano la
superficie laterale e tutte si richiudono su se stesse. L'unica discontinuit\`a
di $\vec B$ si ha per la componente tangenziale attraverso la superficie
laterale del cilindro.

Le equazioni per \ $\vec H\equiv\vec B-4\pi\vec M$ \ sono
\begin{equation}
  \left\{
  \begin{aligned}
    &\dive\vec H=-4\pi\dive\vec M\\
    &\rot\vec H=0
  \end{aligned}
  \right.
  \tag{15.40}
\end{equation}
da cui si vede che in questo caso (cio\`e in assenza di correnti esterne)
$\vec H$ \`e, come $\vec E$, un campo irrotazionale in tutto lo spazio, quindi
ammette un potenziale: \ $\vec H=-\grad\varphi_m$, e dalla 1$^\mathrm{a}$ delle
(15.40) si ha
\begin{equation}
  \Delta\,\varphi_m=4\pi\dive\vec M
  \tag{15.41}
\end{equation}
quindi le ``cariche magnetiche'' sorgenti di $\vec H$ sono date da
\ $\rho_m=-\dive\vec M$.\\
In effetti, siccome $\vec M$ \`e discontinuo attraverso la superficie del
cilindro ($\vec M=0$ fuori), la (15.41) ha significato solo se pensiamo che
$\vec M$ decresca rapidamente, ma con regolarit\`a, dal valore che ha
all'interno del cilindro al valore nullo fuori, altrimenti la (15.41), valida
separatamente all'interno e all'esterno del cilindro, deve essere completata
dalla
condizione (derivante da $\dive\vec B=0$) che la componente normale di $\vec H$
($\hat n$ verso l'esterno) subisce una discontinuit\`a uguale a $4\pi M_n$,
corrispondente ad una distribuzione superficiale di ``cariche magnetiche''
$\sigma_m=M_n$.\\
Nel caso del cilindro uniformemente magnetizzato abbiamo una distribuzione
superficiale di sorgenti (per $\vec H$) sulle basi del cilindro, positive sulla
base verso cui \`e diretto $\vec M$ (il ``polo nord'' della calamita), negative
sull'altra (``polo sud''). Quindi il campo $\vec H$ \`e identico, sia dentro il
cilindro che fuori, al campo elettrico di un cilindro uniformemente polarizzato
o, equivalentemente, di due distribuzioni superficiali di carica $\sigma$ e
$-\sigma$ sulle due basi del cilindro. In particolare, dentro il cilindro
$\vec H$ non \`e uniforme e le sue linee di forza vanno in verso opposto a
quelle di $\vec M$ (dal ``polo nord'' al ``polo sud'') e sono incurvate verso
l'interno, come quelle del campo elettrico dentro un condensatore nel quale le
dimensioni delle armature sono confrontabili con la loro distanza (gli ``effetti
di bordo''). La componente normale di $\vec H$ subisce una discontinuit\`a in
modulo uguale a $4\pi M$ attraverso le basi del cilindro e all'esterno $\vec B$
e $\vec H$ coincidono.

\paragraph{Esempi}
\begin{enumerate}
  \item Consideriamo due magneti, uno a forma di cilindro lungo e sottile, e
    l'altro a forma di cilindro largo e basso (figura 15.2). Supponiamo che la
    magnetizzazione, diretta lungo l'asse, sia la stessa nei due cilindri\\
    Nel primo caso (cilindro lungo e stretto) possiamo approssimare il campo
    magnetico $\vec B$ dentro al cilindro con il campo di un solenoide infinito
    nel quale $nI\equiv\mathcal F=cM$, quindi (all'interno e lontano dalle basi)
    \fig{fig15_2}{15.2}
    \begin{equation}
      B_z=\frac{4\pi In}{c}=4\pi M;\qquad\vec H=\vec B-4\pi\vec M=0.
      \tag{15.42}
    \end{equation}
    Oppure possiamo considerare il magnete come due distribuzioni superficiali
    di cariche magnetiche $\sigma_m=M$ e $\sigma_m=-M$ sulle due basi; se $a$
    \`e il raggio del cilindro e $2L$ la sua lunghezza, siccome $a\ll L$, a
    grande distanza dalle basi il campo $\vec H$ \`e quello di due cariche
    magnetiche puntiformi \ $\pm q_m=\pm\pi a^2\sigma_m=\pm\pi a^2M$; cos\`i,
    per esempio, nel centro del magnete si ha
    \begin{equation}
      H_z=-\frac{2\pi a^2M}{L^2};\qquad B_z=4\pi M-\frac{2\pi a^2M}{L^2}\,\cdot
      \tag{15.43}
    \end{equation}
    Non c'\`e contraddizione fra la (15.42) e la (15.43) in quanto il 2$^\circ$
    metodo fornisce, in prossimit\`a del centro, la correzione di ordine
    $(a/L)^2$ all'approssimazione di solenoide infinito e la 2$^\mathrm{a}$
    delle (15.43) \`e, allo stesso ordine, in accordo con la (7.28).\\
    Con il 1$^\circ$ metodo, quello della corrente superficiale, possiamo
    concludere che $B_z$ sull'asse e subito fuori del magnete vale con buona
    approssimazione $2\pi M$ (la met\`a del valore che ha nel centro) e quindi
    (o anche con il metodo delle cariche magnetiche) $H_z$ vicino alle basi e
    all'interno del magnete vale $-2\pi M$.\\
    Consideriamo ora il secondo magnete. Se $d$ \`e la distanza fra le basi (di
    raggio $b$), siccome $d\ll b$, all'interno e lontano dai bordi $\vec H$ \`e
    come il campo elettrico dentro a un condensatore a facce piane parallele
    infinite, cio\`e
    \begin{equation}
      H_z=-4\pi\sigma_m=-4\pi M;\qquad\vec B=\vec H+4\pi\vec M=0
      \tag{15.44}
    \end{equation}
    ed in questa approssimazione (condensatore ``infinito'') all'esterno del
    magnete $\vec H$ (e quindi anche $\vec B$) \`e nullo. In alternativa, sempre
    grazie al fatto che $d\ll b$, possiamo considerare il magnete come una spira
    di raggio $b$ percorsa dalla corrente $I=cMd$, quindi il campo magnetico
    $\vec B$ dentro e fuori, ma vicino al centro del magnete, per la (7.26) \`e
    \begin{equation}
      B_z=\frac{2\pi Md}{b};\qquad H_z^{\mathrm i}=-4\pi M+\frac{2\pi Md}{b}
      \tag{15.45}
    \end{equation}
    quindi in questo caso il metodo della corrente fornisce la correzione di
    ordine $d/b$ all'approssimazione del condensatore infinito.\\
    Per entrambi i magneti il momento magnetico totale \`e
    $\vec m=\vec M\cdot V$, dove $V$ \`e il volume del magnete, quindi se i due
    magneti hanno lo stesso volume generano \underline{a grande distanza} lo
    stesso campo (campo di dipolo); ovviamente a grande distanza ($r\gg L$,
    \ $r\gg b$) le approssimazioni di solenoide infinito o di condensatore
    infinito non sono valide.

  \item Consideriamo ora un magnete toroidale: per esempio il magnete lungo e
    stretto della figura 15.2 richiuso su se stesso in forma di anello (figura
    15.3). Supponiamo che la magnetizzazione abbia solo la componente
    tangenziale $M_\theta$ (v.\ figura 7.6), costante.
    \fig{fig15_3}{15.3}
    Siccome $\dive\vec M=0$ (perch\'e il flusso di $\vec M$ uscente da qualsiasi
    volume \`e nullo) e sulla superficie $M_n=0$, il campo $\vec H$ non ha
    sorgenti, quindi \`e ovunque nullo. Quindi dentro al magnete c'\`e solo il
    campo magnetico $\vec B=4\pi\vec M$ e all'esterno $\vec B=\vec H=0$. Il
    risultato \`e indipendente dalla forma della sezione del toro.\\
    Se il magnete \`e aperto, come nella figura 15.4, ma la distanza fra le due
    superfici esterne \`e piccola rispetto alle loro dimensioni, il campo
    magnetico $\vec B$ all'interno \`e praticamente uguale a quello del magnete
    della figura 15.3 e siccome $\vec B$ \`e continuo attraverso i ``poli'' del
    magnete, anche fra i poli vale (circa) $4\pi M$, come $H$.
    \fig{fig15_4}{15.4}
\end{enumerate}

\`E importante a questo punto farsi un'idea di quanto pu\`o valere la
magnetizzazione $M$ in un magnete permanente. Partiamo dall'osservazione che i
materiali ferromagnetici diventano paramagnetici al disopra della temperatura
$\Theta$ di Curie, quindi si tratta di sostanze i cui atomi dispongono di un
momento magnetico proprio; in secondo luogo sappiamo che una sostanza
ferromagnetica in un campo magnetico non uniforme risente di una forza
proporzionale al campo, e non al quadrato del campo come le sostanze
diamagnetiche e paramagnetiche. Nel caso di queste ultime, come nel caso dei
dielettrici, la proporzionalit\`a al quadrato del campo deriva dal fatto che \`e
il campo stesso che produce il momento magnetico sul quale poi esercita la
forza. Questo fa supporre che nelle sostanze ferromagnetiche i momenti magnetici
atomici risultino gi\`a quasi tutti allineati nella stessa direzione, nonostante
l'effetto disallineante della temperatura. Il motivo per cui questo succede si
\`e cominciato a capire solo intorno al 1930 (ma tutt'ora
il problema non \`e completamente risolto), in quanto il meccanismo che
favorisce cos\`i fortemente l'allineamento dei momenti magnetici non \`e
attribuibile alla forza che si esercita fra due magneti atomici (che anzi,
quando questi sono affiancati, tenderebbe a disporli antiparalleli), ma \`e di
natura quantistica e ha le sue radici nel principio di esclusione di Pauli.\\
Consideriamo il ferro magnetizzato. Il momento magnetico proprio di ogni atomo
di ferro \`e uguale a due magnetoni di Bohr, il numero atomico del ferro \`e
(circa) 56 (in 56 grammi di ferro ci sono $N_A=6\times10^{23}$ atomi) e la
densit\`a \`e $7.9\,\mathrm{g/cm}^3$, quindi la densit\`a di magnetizzazione di
saturazione, cio\`e se tutti i momenti magnetici sono allineati, \`e
\begin{equation}
  M=\frac{7.9}{56}\cdot6\times10^{23}\cdot2\cdot0.9\times10^{-20}
  \simeq1500\,\mathrm{gauss}
  \tag{15.46}
\end{equation}
quindi il campo magnetico del ferro dovuto alla magnetizzazione \`e dell'ordine
di
\begin{equation}
  4\pi M\approx2\times10^4\,\mathrm{gauss}
  \tag{15.47}
\end{equation}
in accordo con l'esperienza. Siccome abbiamo visto che nelle sostanze
diamagnetiche il momento magnetico indotto \`e dell'ordine di
$10^{-5}\div10^{-6}$ magnetoni di Bohr $\mu_{\textsc b}$, questo \`e anche
l'ordine di grandezza del rapporto fra la magnetizzazione di una sostanza
diamagnetica (o anche paramagnetica a temperatura ambiente) e una
ferromagnetica.\\
Una magnetizzazione come quella data dalla (15.46) corrisponde ad una densit\`a
superficiale di corrente enorme:
\begin{equation}
  \mathcal F=cM=4.5\times10^{13}\,\mathrm{ues}=1.5\times10^4\,\mathrm{A/cm}.
  \tag{15.48}
\end{equation}
Il campo magnetico dato dalla (15.47) \`e quello di saturazione, cio\`e
praticamente il massimo che si pu\`o ottenere all'interno del ferro
magnetizzato; il campo all'interno di una comune calamita, a parte il fatto che
come abbiamo visto negli esempi precedenti dipende (anche) dalla sua forma, \`e
dell'ordine di $10^4\,$gauss, pur sempre un valore ragguardevole. Inoltre, se
prendiamo un pezzo di ferro, a meno che questo non sia stato preventivamente
trattato in maniera opportuna (c'\`e ferro e ferro), non manifesta alcuna
magnetizzazione (invece le forbici di acciaio, cio\`e ferro temprato, sono quasi
sempre magnetizzate).\\
Ma allora, nel ferro comune i momenti magnetici atomici sono allineati o no? Un
pezzo di ferro, come le altre sostanze ferromagnetiche, \`e costituito da
regioni, chiamate \textit{domini di Weiss}, nelle quali i momenti magnetici
atomici sono effettivamente allineati (lungo particolari direzioni determinate
dalla struttura cristallina del ferro), ma nel passaggio da un dominio di Weiss
all'altro la direzione dell'allineamento cambia, ed il risultato sul piano
macroscopico \`e l'assenza di magnetizzazione. Le dimensioni di queste regioni
sono dell'ordine di qualche micron e possono anche essere viste al
microscopio.\\
Evidentemente, per qualche ragione (che ha a che fare con una sorta di
competizione fra l'energia di superficie, quella che separa un dominio
dall'altro, con l'energia di volume, quella associata al volume dei domini), la
situazione in cui nei diversi domini i momenti magnetici atomici non sono
allineati nella stessa direzione \`e energeticamente quella pi\`u favorevole;
questa situazione pu\`o per\`o essere modificata mediante l'applicazione di un
campo magnetico esterno, il cui effetto \`e quello di rendere energeticamente
pi\`u favorevoli domini sempre pi\`u grandi in cui i momenti magnetici atomici
sono allineati con il campo esterno: \`e addirittura possibile registrare
sperimentalmente il ``ru\-
more magnetico'', opportunamente amplificato e trasformato in rumore acustico,
che si produce quando un intero dominio (che contiene diversi miliardi di atomi)
cambia bruscamente la propria magnetizzazione per allinearsi con i domini
vicini, gi\`a allineati con il campo (\textit{effetto Barkhausen}). Il
risultato, aumentando il campo esterno, \`e un progressivo estendersi delle
dimensioni dei domini ``ben allineati'' a spese degli altri, e quindi la
magnetizzazione del nostro campione.

Il modo in cui i diversi materiali ferromagnetici si magnetizzano quando viene
applicato un campo esterno dipende dal tipo di materiale, dal trattamento
meccanico subito (p.es.\ la tempratura), la presenza di impurit\`a, ecc. Senza
addentrarci nella enorme casistica delle possibili situazioni, consideriamo due
casi abbastanza emblematici, quello del ferro dolce, e quello dell'acciaio (che
sono due rappresentanti della categoria dei materiali ``magneticamente dolci'' e
di quelli ``magneticamente rigidi'').

Supponiamo di avere due avvolgimenti toroidali sovrapposti, come quelli
considerati nell'esempio 1 nel capitolo 8.7, entrambi avvolti su un supporto
ferromagnetico sagomato in forma toroidale; nel primario, mediante un generatore
di corrente, facciamo passare una corrente $I(t)$ inizialmente nulla e che
aumentiamo gradualmente. Mediante un voltmetro misuriamo la f.e.m.\ indotta nel
secondario; nella figura 15.5 \`e riportata la rappresentazione circuitale di
questo sistema: le righe nere fra le due induttanze significano che primario e
secondario sono avvolti sullo stesso supporto magnetico e la freccia attraverso
il generatore di corrente significa che \`e possibile variare la corrente
erogata.
\fig{fig15_5}{15.5}
La conoscenza della corrente comporta, grazie all'ultima delle (15.9) e al
teorema della circuitazione, la conoscenza di $H$ nel supporto ferromagnetico:
la (7.45) ora vale con $H$ al posto di $B$
\begin{equation}
  H_\theta(t)=\frac{2NI(t)}{c\varrho}\,\cdot
  \tag{15.49}
\end{equation}
Se il supporto toroidale \`e sottile possiamo considerare $H_\theta$ costante
sulla sua sezione.

Variando la corrente nel primario si ha una f.e.m.\ indotta nel secondario che,
siccome la 3$^\mathrm{a}$ delle (15.9) \`e uguale all'equazione nel vuoto, \`e
sempre data dalla (8.2), e quindi \`e proporzionale a $\dot B$ (e non ad
$\dot H$); la conoscenza di $\dot B(t)$ in ultima analisi permette di
determinare $B$ ad ogni istante, e quindi $B$ in funzione di $H$; se il
comportamento dei materiali ferromagnetici fosse lineare la curva $B=B(H)$
sarebbe una retta, e il coefficiente angolare di questa retta sarebbe la
permeabilit\`a magnetica $\mu$; sperimentalmente invece si ottengono curve del
tipo di quella riportata nella figura 15.6 (\textit{curva di magnetizzazione}),
qualunque sia la sostanza ferromagnetica (ferro dolce, acciaio, \dots).
\fig{fig15_6}{15.6}
Va notato che la scala sull'ordinata e quella sull'ascissa sono ben diverse,
circa per un fattore $10^3$: il campo magnetico di saturazione
$B_{\mathrm{sat}}$, quello che nella figura 15.6 corrisponde alla linea
tratteggiata orizzontale, nel caso del ferro dolce e dell'acciaio \`e un campo
dell'ordine di $15000\,$gauss, che si raggiunge quando il campo $H$ vale circa
$30\,$oersted ($1\,$oersted $=1\,$gauss!) nel caso del ferro dolce, e (circa)
$200\,$oersted nel caso dell'acciaio. Quindi possiamo dire che praticamente
tutto il campo $\vec B_{\mathrm{sat}}$ \`e dovuto alla magnetizzazione:
\ $\vec B_{\mathrm{sat}}\simeq4\pi\vec M_{\mathrm{sat}}$.\\
Per la (15.49), se $\varrho=10\,$cm e $N=1000$ (possiamo realizzare il primario
con un avvolgimento a pi\`u strati), si ottiene un campo $H$ di $200\,$oersted
con una corrente $I$ di $10\,$A che, siccome le spire sono 1000 su una
circonferenza di $\simeq60\,$cm, corrisponde ad una corrente superficiale di
circa $170\,$A/cm: una discreta corrente, ma irrisoria rispetto a quella di
magnetizzazione \big(v.\ (15.48)\big); la (relativamente) piccola corrente $I$
\`e servita solo a ``mettere in ordine'' i domini di Weiss.

Fino a questo punto il comportamento del ferro dolce e dell'acciaio non \`e
sostanzialmente diverso: l'acciaio \`e solo un po' pi\`u ``rigido'' nel
raggiungere la saturazione.\\
La vera differenza si manifesta se ora, dopo avere raggiunto $B_{\mathrm{sat}}$,
diminuiamo la corrente $I$, e quindi il campo $H$. Nel caso dei materiali
``magneticamente (molto) dolci'' la curva di magnetizzazione della figura 15.6
viene percorsa all'inverso (in realt\`a solo approssimativamente), e quindi il
processo di magnetizzazione \`e reversibile o, detto in altri termini, $\vec B$
\`e una funzione univoca di $\vec H$ e questa relazione funzionale \`e
l'equazione costitutiva che bisogna aggiungere alle (15.9). In questi casi
(materiali magneticamente dolci), almeno quando per piccoli valori di $H$ la
curva $B(H)$ pu\`o essere approssimata con una retta, \`e possibile definire la
permeabilit\`a magnetica del mezzo dalla relazione
\begin{equation}
  \vec B=\mu\vec H
  \tag{15.50}
\end{equation}
e, come \`e gi\`a abbastanza evidente dal grafico 15.6, i valori di $\mu$ sono
molto grandi: possono variare da $\mu\simeq10$ a $\mu\simeq10^4$ e in casi
speciali pu\`o addirittura arrivare a valori $\mu\simeq10^6$.

Invece, nel caso dei materiali ``magneticamente rigidi'', se dopo avere
raggiunto la saturazione (lungo la linea tratteggiata nella figura 15.7)
diminuiamo il campo $H$, $B$ segue la curva indicata dalla freccia verso
sinistra, e quindi diminuisce molto lentamente e non \`e nullo quando $H$
(cio\`e la corrente nel primario) \`e nulla: questa \`e il \textit{campo
residuo} nel nostro campione, dovuto alla \textit{magnetizzazione residua}: il
campione \`e quindi \`e diventato un magnete permanente, con un campo interno
non molto inferiore al campo di saturazione, cio\`e circa $10^4\,$gauss. Per
smagnetizzare il magnete \`e necessario invertire la corrente in modo da avere
un campo $H$ in verso opposto: il valore $H_c$ che corrisponde nel grafico 15.7
a $B=0$ \`e detto \textit{forza coercitiva} e vale, nel caso dell'acciaio, circa
$60\,$oersted.
\fig{fig15_7}{15.7}
La curva a tratto pieno nella figura 15.7, che si ottiene variando con
continuit\`a la corrente nel primario, \`e detta \textbf{ciclo di isteresi}; se
poi, in un punto del ciclo di isteresi dove la pendenza \`e grande cerchiamo di
invertire il processo (se la corrente stava diminuendo ora la facciamo
aumentare), $B$ non ripercorre all'inverso la curva, ma segue un percorso
diverso: il ciclo di isteresi non \`e reversibile (torneremo su questo punto nel
paragrafo 7).\\
Quindi $B$ non \`e una funzione univoca di $H$, in quanto dipende dai valori
precedentemente assunti da $H$, cio\`e dalla storia passata. In questo caso non
\`e possibile scrivere
un'equazione costitutiva che lega fra loro $B$ e $H$.

La ricerca per la realizzazione di materiali con ben precise caratteristiche
magnetiche \`e molto attiva, soprattutto sotto la spinta di uno sviluppo
tecnologico che fa sempre maggior uso di questi materiali: fra i materiali
``dolci'', necessari per la costruzione dei motori e dei trasformatori, citiamo
i \textit{permalloy}, sostanzialmente leghe di nichel e ferro, che hanno un
elevato campo magnetico residuo ($\simeq5000\,$gauss) e una piccolissima forza
coercitiva ($\simeq4\times10^{-3}\,$oersted); fra i materiali ``rigidi'',
necessari per le memorie magnetiche, i magneti permanenti, ecc., il citatissimo
Alnico V, una lega che contiene in percentuali ben precise oltre all'alluminio,
il nichel e il cobalto (da cui il nome), anche ferro e rame, che ha un campo
residuo di $13000\,$gauss e una forza coercitiva di $550\,$oersted.

\subsection{Schermaggio magnetico}

Se vogliamo schermarci dal campo elettrico, ci mettiamo dentro la cavit\`a di un
conduttore (gabbia di Faraday): l'effetto schermante deriva dal fatto che in un
conduttore il campo elettrico \`e nullo -- quindi il conduttore \`e una regione
equipotenziale -- e dall'unicit\`a della soluzione dell'equazione di Laplace,
date le condizioni al contorno (problemi di Dirichlet e di Neumann).\\
Esiste la possibilit\`a di schermarci da un campo magnetico esterno, per esempio
dal campo magnetico terrestre? Il problema \`e di notevole rilevanza pratica in
quanto in diverse situazioni (funzionamento di apparecchiature elettroniche,
misure in assenza di campi magnetici, \dots) \`e necessario creare regioni nelle
quali la presenza, e quindi gli effetti, di campi magnetici esterni siano
praticamente nulli. Vedremo che i materiali magneticamente dolci di elevata
permeabilit\`a magnetica ($\mu\to\infty$) si comportano, per quanto riguarda
$\vec H$, come i conduttori nei confronti del campo elettrico: uno dei materiali
particolarmente adatto a questo scopo \`e il \textit{Supermalloy} per il quale,
per piccoli campi, si pu\`o scrivere la (15.50), con $\mu\simeq10^6$.\\
L'idea \`e questa: siccome $B$ all'interno di un materiale magnetico (dolce) \`e
dell'ordine del campo esterno $B_0$, ed $H=B/\mu$, se $\mu\to\infty$ ne segue
che $H\to0$; se dentro il materiale $H=0$ siamo nella stessa situazione di un
conduttore nei confronti del campo elettrico.\\
Il punto pi\`u delicato \`e che dobbiamo convincerci che effettivamente dentro a
un materiale di elevata permeabilit\`a magnetica (e quindi facilmente
magnetizzabile) $B$ ``resta finito'', cio\`e dell'ordine di $B_0$.\\
Consideriamo per semplicit\`a una sfera del nostro materiale in un campo esterno
$\vec B_0\equiv\vec H_0$ uniforme: siccome le equazioni per $\vec H$ sono
identiche a quelle del campo elettrico (in una regione priva di correnti di
conduzione), il problema di determinare $\vec H$ all'interno della sfera
magnetizzabile \`e identico a quello di determinare il campo elettrico dentro
una sfera di dielettrico in un campo esterno, problema la cui soluzione \`e data
dalla (14.47); quindi dentro la sfera magnetizzabile
\begin{equation}
  \vec H=\frac{3}{2+\mu}\vec H_0
  \quad\Rightarrow\quad\vec B=\frac{3\mu}{2+\mu}\vec B_0
  \tag{15.51}
\end{equation}
da cui si vede che effettivamente se $\mu\to\infty$, \ $H\to0$,
\ $B\to\,\approx B_0$.\\
Supponiamo di mettere un corpo cavo di materiale magneticamente dolce con ele\-
vata permeabilit\`a magnetica in un campo esterno $\vec B_0$ (per esempio, una
sfera cava). Nell'approssimazione $\mu\to\infty$ tutta la regione occupata dal
materiale \`e una regione equipotenziale per il potenziale magnetico $\varphi_m$
(perch\'e $H=0$), e quindi \`e equipotenziale anche la cavit\`a (unicit\`a della
soluzione dell'equazione $\Delta\,\varphi_m=0$). Ma dentro la cavit\`a, cio\`e
nel vuoto, \ $\vec B=\vec H\ \Rightarrow\ B=0$ \ quindi, nel limite
$\mu\to\infty$, abbiamo realizzato lo schermo magnetico.\\
Nella regione occupata dal materiale (fra l'esterno e la cavit\`a) il campo
magnetico $\vec B$ non \`e nullo e come conseguenza del fatto che \`e un campo
solenoidale ($\dive\vec B=0$) le sue linee di forza hanno un andamento molto
particolare: siccome dentro $\vec H=0$, fuori $\vec H=\vec B$ e la componente
tangenziale di $\vec H$ \`e continua attraverso la superficie, ne segue che
$\vec H$, e quindi $\vec B$, \`e ortogonale alla superficie esterna del
materiale (come il campo elettrico sulla superficie di un conduttore); entrando,
la componente tangenziale di $\vec B$ subisce una discontinuit\`a (a causa della
corrente superficiale di magnetizzazione) che ha l'effetto di allontanare le
linee di forza dalla normale e incanalarle nella regione fra la superficie
esterna e quella interna; siccome la componente normale di $\vec B$ \`e continua
e nella cavit\`a $\vec B=0$, segue che le linee di forza di $\vec B$ che
arrivano sulla superficie interna del materiale sono tangenti ad essa. Ne segue,
specialmente se la regione occupata dal materiale \`e sottile (p.es.\ una corona
sferica di raggi interno ed esterno quasi uguali) che tutte le linee di forza
che entrano nel materiale si addensano nel piccolo volume a disposizione: questa
concentrazione delle linee di forza nei materiali magnetici ad alta
permeabilit\`a permette di aumentare notevolmente i coefficienti di auto e mutua
induzione di avvolgimenti realizzati attorno a supporti magnetici (chiamati
\textit{nuclei magnetici}), in quanto la presenza del nucleo magnetico provoca
un aumento dell'intensit\`a del campo e una riduzione del \textit{flusso
disperso} (cio\`e non concatenato dal circuito).

\paragraph{Esempi}
\begin{enumerate}
  \item Supponiamo di mettere una corona sferica di un materiale di
    permeabilit\`a magnetica $\mu$ in un campo esterno uniforme $\vec B_0$.
    Vogliamo stimare l'intensit\`a del campo magnetico all'interno della
    cavit\`a. La soluzione esatta del problema \`e possibile ma complicata per
    cui, se ci accontentiamo di una stima, faremo delle approssimazioni.\\
    Se non ci fosse la cavit\`a, il campo $\vec H_{\mathrm{sfera}}$ dentro la
    sfera sarebbe quello dato dalla (15.51). Se ora supponiamo che il raggio
    esterno della corona sia molto maggiore di quello interno, possiamo
    approssimare il campo $\vec H$ nella cavit\`a come il campo in una cavit\`a
    sferica in un mezzo infinito di permeabilit\`a $\mu$ nel quale \`e presente
    (in assenza della cavit\`a) il campo uniforme $\vec H_{\mathrm{sfera}}$: la
    soluzione, fatte le opportune sostituzioni \`e data dalla (14.49)
    \begin{equation}
      \vec H_{\mathrm{cav}}\equiv\vec B_{\mathrm{cav}}
      \simeq\frac{3\mu}{2\mu+1}\vec H_{\mathrm{sfera}}
      =\frac{9\mu}{(2+\mu)(2\mu+1)}\vec B_0\to\frac{9}{2\mu}\vec B_0
      \tag{15.52}
    \end{equation}
    (l'ultimo passaggio corrisponde al caso $\mu\gg1$).

  \item Ne capitolo 8.4 abbiamo calcolato il coefficiente di autoinduzione di un
    avvolgimento toroidale a spire rettangolari ($N$ spire, raggio interno $a$,
    esterno $b$, altezza $h$). Vediamo come si modifica il risultato, dato dalla
    (8.32), se le spire sono avvolte su un materiale ferromagnetico dolce per il
    quale vale la (15.50).
    La definizione di coefficiente di auto (o mutua) induzione \`e sempre la
    stessa: $\Phi(\vec B)/cI$, solo che ora la corrente $I$ determina $\vec H$
    tramite l'ultima delle (15.9) \ $\rot\vec H=4\pi\vec J/c$, ed \`e dato dalla
    (15.49): la stessa espressione che si ha nel vuoto per $\vec B$
    \big(v.\ (7.45)\big). Siccome assumiamo che valga la (15.50) si ha (se il
    materiale \`e omogeneo)
    \begin{equation}
      \Phi(\vec B)=\Phi(\mu\vec H)=\mu\Phi(\vec H)
      \tag{15.53}
    \end{equation}
    quindi il coefficiente di autoinduzione risulta $\mu$ volte quello nel
    vuoto:
    \begin{equation}
      L=\frac{2\mu N^2h}{c^2}\log\frac{b}{a}\,\cdot
      \tag{15.54}
    \end{equation}
\end{enumerate}

\subsection{L'energia magnetica}

Come nel caso dei dielettrici, si arriva al concetto di energia del sistema
campo + materia a partire dal lavoro che occorre fare sul sistema per instaurare
il campo magnetico. Siccome nel caso dei dielettrici ci siamo limitati al caso
in cui \`e valida l'approssimazione lineare, non sempre le conclusioni raggiunte
in quel caso possono essere trasferite al caso magnetico, dato che le sostanze
magneticamente interessanti di solito sono quelle ferromagnetiche.\\
Come abbiamo gi\`a discusso nel caso dei dielettrici, le forze esterne sono
quelle che i generatori devono esercitare contro le forze del campo sulle
cariche ``vere'' per instaurare i campi: il lavoro fatto nell'unit\`a di tempo
(potenza) dalle forze esterne \`e dato da
\begin{equation}
  W^\mathrm{est}(t)=-\int\vec E\cdot\vec J\,dV.
  \tag{15.55}
\end{equation}
La (15.55) formalmente coincide con $-W(t)$ dato dalla (11.3) (nella (11.3) $W$
\`e la potenza delle forze del campo), ma il contenuto \`e diverso. Nella
(15.55) $\vec J$ \`e solo la corrente di conduzione (quella dovuta al moto delle
cariche ``vere''), e non contiene le correnti microscopiche, quelle che
compaiono nelle equazioni (14.4) per i campi microscopici e che, mediate, danno
luogo alle correnti di polarizzazione e a quelle di magnetizzazione
\big(equazione (15.5)\big). Nella (11.3), invece, $\vec J$ \`e tutta la
corrente, sia quella di conduzione che le correnti microscopiche. Come abbiamo
gi\`a rilevato a proposito della discussione fatta nel caso dei dielettrici, la
(11.3) in pratica ha senso quando ha senso fare una chiara distinzione fra campo
e materia (elettromagnetismo nel vuoto).

Per determinare in generale quella che normalmente viene chiamata energia del
sistema campo + materia
\begin{equation}
  U(t)=\int_{-\infty}^tW^\mathrm{est}(t')\,dt'
  \tag{15.56}
\end{equation}
(ma che pi\`u correttamente fino a prova contraria dovremmo limitarci a chiamare
``lavoro fatto sul sistema''), si procede come nel capitolo 11.1. Occorre per\`o
usare le equazioni (15.9), dato che in esse $\vec J$ \`e la corrente di
conduzione: moltiplichiamo scalarmente la 4$^\mathrm{a}$ e la 2$^\mathrm{a}$
delle (15.9) rispettivamente per $\vec E$ e per $\vec H$ e sottraiamole:
\begin{equation}
  \vec E\cdot\rot\vec H-\vec H\cdot\rot\vec E=\frac{4\pi}{c}\vec E\cdot\vec J
  +\frac1c\Big(\vec E\cdot\frac{\partial\vec D}{\partial t}
  +\vec H\cdot\frac{\partial\vec B}{\partial t}\Big)
  \tag{15.57}
\end{equation}
quindi grazie all'identit\`a (11.6) si ottiene
\begin{equation}
  -\vec E\cdot\vec J
  =\frac{1}{4\pi}\Big(\vec E\cdot\frac{\partial\vec D}{\partial t}
  +\vec H\cdot\frac{\partial\vec B}{\partial t}\Big)
  +\frac{c}{4\pi}\dive(\vec E\wedge\vec H)
  \tag{15.58}
\end{equation}
e quindi
\begin{equation}
  W^\mathrm{est}(t)
  =\frac{1}{4\pi}\int\Big(\vec E\cdot\frac{\partial\vec D}{\partial t}
  +\vec H\cdot\frac{\partial\vec B}{\partial t}\Big)\,dV.
  \tag{15.59}
\end{equation}
La (15.59), come le (15.9), \`e indipendente dalle equazioni costitutive che
legano $\vec E$, $\vec B$ a $\vec D$ e $\vec H$ e anche dal fatto che queste
esistano o meno; in generale non \`e possibile ricavare, integrando su $t$,
l'energia del sistema $U(t)$ in termini dei campi all'istante $t$, perch\'e
abbiamo visto che ci sono casi in cui i campi ad un dato istante, e quindi anche
l'energia del sistema, dipendono dalla storia precedente del sistema (fenomeno
dell'isteresi). Se il comportamento del mezzo \`e lineare, allora
$\vec B=\mu\vec H$, \ $\vec D=\epsilon\vec E$, per cui
\begin{equation}
  \int\vec E\cdot\frac{\partial\vec D}{\partial t'}\,dt'
  =\frac12\vec E\cdot\vec D\,,
  \qquad\int\vec H\cdot\frac{\partial\vec B}{\partial t'}\,dt'
  =\frac12\vec H\cdot\vec B
  \tag{15.60}
\end{equation}
e in questo caso quindi
\begin{equation}
  U(t)=\frac{1}{8\pi}\int\big(\vec E\cdot\vec D+\vec H\cdot\vec B\big)\,dV.
  \tag{15.61}
\end{equation}
La (15.61) si riduce alla (14.57) nel caso che sia presente solo un campo
elettrico statico, e nel caso che sia presente solo un campo magnetico statico
d\`a
\begin{equation}
  U_M=\frac{1}{8\pi}\int\vec H\cdot\vec B\,dV.
  \tag{15.62}
\end{equation}

\paragraph{Esempio}
\begin{quote}
  Riprendiamo in considerazione l'avvolgimento toroidale dell'esempio 2 del
  paragrafo precedente e calcoliamo il lavoro che un generatore deve fare per
  far circolare una corrente $I$ nel circuito.\\
  Si ha
  \begin{equation}
    \begin{split}
      &\frac{dU_M(t)}{dt}=W^\mathrm{est}(t)=-\mathcal E(t)I(t)=L\dot I(t)I(t)
      \quad\Rightarrow\\
      &U_M=\frac12LI^2=\frac{I^2}{2}\frac{2\mu N^2h}{c^2}\log\frac ba\,\cdot
    \end{split}
    \tag{15.63}
  \end{equation}
  Se invece usiamo la (15.62), per la (15.49) si ha
  \[
    U_M=\frac{\mu}{8\pi}\int H^2\,dV
    =\frac{\mu h}{8\pi}\int_0^{2\pi}d\phi\int_a^bH^2\varrho\,d\varrho
    =\frac{\mu hN^2I^2}{c^2}\log\frac ba
  \]
  che coincide con la (15.63).
\end{quote}

Nel capitolo 14.6 abbiamo messo in evidenza che la (14.57) rappresenta, per un
dielettrico mantenuto a temperatura costante e per trasformazioni reversibili,
la variazione $\Delta\,\mathcal F$ della sua energia libera
$\mathcal F=\mathcal U-T\mathcal S$ tra lo stato con il campo e quello senza
campo. Il problema del significato dell'energia magnetica $U_M$ si pone negli
stessi termini e nel
caso delle sostanze diamagnetiche o paramagnetiche le conclusioni sono le stesse
come nel caso dei dielettrici, per\`o in questi casi sappiamo che gli effetti
della presenza della materia sul campo magnetico, e quindi sull'energia, sono di
solito molto piccoli.

Il caso importante \`e quello delle sostanze ferromagnetiche, anche perch\'e in
questi casi le trasformazioni possono essere irreversibili.\\
Supponiamo di fare effettuare ad una sostanza ferromagnetica il ciclo di
isteresi a contatto con l'ambiente, quindi in condizioni isoterme,
nell'intervallo di tempo $t_1$, $t_2$. Il lavoro fatto sul sistema per la
(15.59) \`e dato da
\begin{equation}
  L^\mathrm{est}=\frac{1}{4\pi}\int_{t_1}^{t_2}dt
  \Big(\int\vec H\cdot\frac{\partial\vec B}{\partial t}\,dV\Big).
  \tag{15.64}
\end{equation}
Supponiamo, ma solo per semplicit\`a, che $B$ ed $H$ siano costanti dentro il
materiale, allora siccome nell'intervallo di tempo $(t_1,\,t_2)$ \ $\vec B$ ed
$\vec H$ compiono un ciclo, se $V$ \`e il volume del campione, possiamo
riscrivere la (15.64) come
\begin{equation}
  L^\mathrm{est}=\frac{V}{4\pi}\oint\vec H\cdot d\vec B
  \tag{15.65}
\end{equation}
da cui si vede che, se $\vec B$ e $\vec H$ sono paralleli (come nel caso di un
magnete toroidale) $L^\mathrm{est}$ \`e dato dall'area racchiusa dal ciclo di
isteresi nel piano $(H,\,B)$, proprio come nel piano di Clapeyron $(V,\,p)$ il
lavoro fatto sul sistema in un ciclo \`e l'area racchiusa dal ciclo ($L>0$ se il
ciclo \`e percorso in senso antiorario). Poich\'e $L^\mathrm{est}\neq0$ e il
ciclo \`e monotermo, necessariamente \`e irreversibile e $L^\mathrm{est}>0$
(postulato di Lord Kelvin).\\
Siccome il sistema ha compiuto un ciclo, la sua energia libera (e anche la sua
energia interna) non \`e cambiata, quindi il lavoro fatto eguaglia interamente
il calore ceduto all'ambiente (il nucleo ferromagnetico si scalda e mediante
l'opportuno sistema di raffreddamento il calore viene trasferito all'esterno).\\
Quindi in questi casi (trasformazioni irreversibili) la (15.59) non rappresenta
la variazione di qualche forma di energia (interna, libera, \dots), ma solo il
lavoro fatto nell'unit\`a di tempo sul sistema: sappiamo infatti che nei
processi isotermi irreversibili la differenza fra il lavoro fatto e la
variazione dell'energia libera \`e positivo ed \`e uguale a $T\Delta\,S$, dove
l'entropia \`e quella totale del sistema e delle sorgenti; in un ciclo
$T\Delta\,S$ \`e il calore scambiato con l'ambiente.
